% Slide 01 (CH)
\begin{frame}{Il ponte logico-computazionale (CH)}
    \begin{itemize}
        \item \textbf{Intuizione di base}: le proposizioni sono \textbf{tipi}, le prove sono \textbf{programmi}.
        \item Corrispondenza biunivoca tra logica intuizionista e calcoli tipati:
    \end{itemize}

    \vskip 0.30cm
    \begin{table}
        \centering
        \small
        \begin{tabular}{lll}
            \toprule
                \textbf{Logica intuizionista} & \textbf{Teoria dei Tipi}   & \textbf{Termine Computazionale}              \\
            \midrule
                Proposizione $A$              & Tipo $A$                   & Termine $t : A$                              \\
                Implicazione $A \to B$        & Tipo funzione $A \to B$    & Astrazione $\lambda x. t$                    \\
                Congiunzione $A \land B$      & Tipo prodotto $A \times B$ & Coppia ordinata $\langle t_1, t_2 \rangle$   \\
                Disgiunzione $A \lor B$       & Tipo somma $A + B$         & Costruttori $\mathsf{inl}, \mathsf{inr}$     \\
                Assurdo $\bot$                & Tipo vuoto $\mathbf{0}$    & Funzione esuriente (assurdo)                 \\
            \bottomrule
        \end{tabular}
    \end{table}

    \vskip 0.20cm
    \noindent
    \begin{block}{Principio fondamentale}
        Dimostrare un teorema equivale a sintetizzare un programma conforme alla specifica di tipo.
    \end{block}
\end{frame}

% Slide 02 (CH)
\begin{frame}{Isomorfismo statico e dinamico}
    \begin{columns}[t, onlytextwidth]

        \begin{column}{0.48\textwidth}
            \begin{block}{Dimensione statica}
                \begin{minipage}[t][4.0cm][t]{\linewidth}
                    \begin{itemize}
                        \item Formule $\longleftrightarrow$ Tipi
                        \item Derivazioni logiche $\longleftrightarrow$ Termini tipati
                        \item \textbf{Validità}:
                        \vskip -0.15cm
                        \[
                            \Gamma \vdash A \iff \Gamma \vdash t : A
                        \]
                        \vskip -0.15cm
                        \item Il \emph{type checking} decide la validità della dimostrazione.
                    \end{itemize}
                \end{minipage}
            \end{block}
        \end{column}

        \begin{column}{0.48\textwidth}
            \begin{block}{Dimensione dinamica}
                \begin{minipage}[t][4.0cm][t]{\linewidth}
                    \begin{itemize}
                    \item Cut-elimination $\longleftrightarrow$ $\beta$-riduzione
                    \item Semplificazione di lemmi come esecuzione.
                    \item \textbf{Normalizzazione}:
                        \vskip -0.15cm
                        \[
                        \mathcal{D} \rightsquigarrow^* \mathcal{D}_{\text{norm}} \iff t \to_\beta^* v
                        \]
                        \vskip -0.15cm
                    \item Prove canoniche: forme normali (prive di redex).
                    \end{itemize}
                \end{minipage}
            \end{block}
        \end{column}
    \end{columns}
\end{frame}